\section{Introduction}

The \cl{}\cite{ansi:common:lisp}
\texttt{read} function, usually referred to as the \cl{}
\emph{reader}, is an obviously essential part of any \cl{}
implementation.  It is normally used as the first step of the
\texttt{read-eval-print} loop or in the file compiler, in both cases
reading forms that are then passed to the evaluator.  Furthermore,
when a \cl{} system is built (or \emph{bootstrapped}), the reader is
needed early since code has to be read before it can be compiled or
evaluated.  But reading native code for compilation or evaluation is
not the only use for a reader.  Some situations where a reader might
be useful are:

\begin{description}
\item[Cross compilation] When a \cl{} system (the \emph{target}) is
  built from source, a common technique is to compile the source on an
  existing \cl{} system (the \emph{host}).  The host system can be
  (perhaps a different version of) the same implementation as the
  target, or it can be a different implementation altogether.  The
  reader is then executed by the host \cl{} system.  In either case,
  it is undesirable for the reader to intern symbols in host packages.
  The \sbcl{} implementation of \cl{} solves this problem by having
  the source code defined in a package that cannot clash with the host
  packages.  The main disadvantage of this approach is that all source
  code that is part of the target system must be specific to that
  system.  External libraries are excluded if such libraries involve
  packages that may clash with some host packages.  With a highly
  configurable reader, cross compilation can be addressed in a
  different way.
\item[Parsing an editor buffer] As a concrete example of this
  use-case, consider the (bad) \cl{} program in
  Listing~\ref{lis:bad-program} which a user might type into an editor
  or integrated development environment (IDE) for \cl{} code.  The IDE
  could support the user with accurate syntax
  highlighting\footnote{For example, the second occurrence of
    \texttt{defun} should be highlighted as a parameter or variable,
    not as the \texttt{cl:defun} operator.  Ironically, the \LaTeX{}
    package used for typesetting the listing gets this wrong.} and
  notify them about the syntactic errors\footnote{Examples of
    syntactic errors are three package markers in \texttt{cl:::nth}, a
    digit using an invalid radix in \texttt{\#b0102} and a missing
    closing parenthesis at the end of the program.}, semantic
  errors\footnote{Only one branch of the nested read-time conditional
    \texttt{\#+sbcl \#b0102 \#-sbcl \#b0100} is reachable.} and style
  issues\footnote{Toplevel comments are specified to use three
    semicolons.  Closing parentheses should not be placed on separate
    lines.}.

  Existing editors for \cl{} code do not provide these features to the
  full extent because they cannot analyze the buffer contents
  faithfully enough.  Many assumptions are made that may not hold
  true.  For example, reader macros are not taken into account, and
  the analysis does not permit the editor to accurately determine the
  role of symbols in the buffer.

  The only technique that can faithfully analyze a buffer of \cl{}
  code is the \cl{} reader.  But then, again, it is undesirable for
  the reader to intern symbols in any package at all, and the analysis
  must not fail for errors such as incorrect package names used as
  package prefixes, incomplete code, or syntax errors in tokens.
  Consequently, an extended \cl{} reader that is configurable and can
  recover from errors is necessary for this approach but not
  sufficient on its own.  Instead, such a reader has to collect enough
  information about errors and the input text for subsequent stages to
  process.
\end{description}

\begin{listing}
\begin{minted}[fontsize=\small,linenos]{lisp}
;; Comment
(defun get-name (defun)
  #+sbcl (cl:::nth
    #+sbcl #b0102 #-sbcl #b0100
    defun
  )
\end{minted}
\caption{A \cl{} program with multiple errors and style issues that a
  user could type into an editor buffer.}
\label{lis:bad-program}
\end{listing}

\noindent
In this paper, we describe \eclector{}; a \cl{} reader that is
independent of any particular \cl{} implementation.  Furthermore,
\eclector{} is highly configurable, so that it may be configured not
to intern symbols in any host package, and so that syntax errors can
be avoided or recovered from.

%%% Local Variables:
%%% mode: LaTeX
%%% TeX-master: "eclector"
%%% End:
